\chapter{Conclusion and Future Work}

\section{Conclusion}

This thesis presented FireWatch: A Drone-Ready Wildfire Risk Prediction and
Fire/Smoke Detection System. The project was developed as a regional
decision-support prototype for Karabük, Türkiye. Its purpose was to show how
weather-based wildfire risk prediction, visual fire/smoke detection,
persistent alert management, dashboard monitoring, scheduler support, and a
drone-ready layer can be combined into one coherent software system.

The prediction methodology used May--October seasonal Karabük records from
2012 to 2025. Open-Meteo daily weather aggregates were used as accessible
model inputs, while the processed dataset included a processed FWI target
associated with official fire-danger data and used for supervised learning.
The final prediction workflow used a two-stage structure:
Stage 1 applied a \texttt{HistGradientBoostingRegressor} to estimate
continuous FWI, and Stage 2 applied a \texttt{RandomForestClassifier} as a
safety-oriented classifier for days where \( \mathrm{FWI} \geq 35 \).

The visual monitoring part of the project used a custom YOLOv8 detector, with
\texttt{best4.pt} as the active fire/smoke detection model. This detector was
kept separate from the weather-based prediction model. Its role was to provide
visual evidence from camera or drone-ready monitoring sources, while the
prediction model provided daily fire-weather risk awareness.

The system implementation connected these models through a FastAPI backend,
SQLite persistence, a scheduler, and a React/TypeScript dashboard. SQLite made
it possible to preserve prediction history and detection alerts, including
review states such as read/unread and soft deletion. The dashboard provided an
operator-facing interface where risk decisions, monitoring feeds, detection
alerts, system/model information, workflow explanation, and map context could
be reviewed together.

The final system answers the research question by demonstrating that daily
risk prediction, live fire/smoke detection, persistent alert management, and
dashboard-based decision support can be integrated into one regional wildfire
monitoring prototype. At the same time, the project remains realistic about
its scope. FireWatch is a B.Sc. engineering prototype and regional pilot; it
is not an official emergency response system and does not claim full
autonomous production drone deployment.

\section{Future Work}

Future work should first improve the robustness of the visual detection
module. The current prototype uses frame-based YOLOv8 detection, but practical
outdoor use would benefit from multi-frame confirmation, glare and reflection
filtering, bounding-box area filtering, and testing on a larger labelled
fire/smoke validation set. A formal benchmark should also be prepared so that
detection metrics such as mAP, precision, and recall can be reported reliably.

The prediction workflow can also be extended. Future versions could include
additional years, neighbouring regions, fuel-related variables, topography,
land-cover information, and human-activity indicators if reliable data are
available. The high-risk threshold and classifier probability threshold could
also be calibrated using operational feedback. Uncertainty estimation would
make the prediction output more useful by showing not only the predicted risk
value, but also the confidence or uncertainty around that decision.

The drone-ready layer should be developed gradually and under clear safety
constraints. Future work may include supervised waypoint patrol, geofencing,
return-to-home behaviour, battery-aware mission planning, and safer mission
state management. These improvements would also require regulatory
compliance, operator procedures, hardware validation, and production-level
safety testing before any real autonomous patrol could be considered.

A production deployment would require stronger infrastructure than the current
local prototype. SQLite is suitable for demonstration and single-operator use,
but a deployed multi-user system should use a server-based database such as
PostgreSQL. Authentication, user roles, backups, monitoring, logging, and
cloud or on-premise deployment procedures would also be required before the
system could be used in an operational environment.

Data pipeline traceability is another important area for improvement. Future
work should formalize the raw Open-Meteo and CEMS/Copernicus download scripts,
store reproducible data-preparation records, and document dataset versions.
This would make it easier to rebuild the processed dataset, audit modelling
decisions, and compare future model versions fairly.

Finally, field validation is necessary before FireWatch can be considered for
real operational use. The system should be tested with outdoor camera scenes,
different lighting conditions, different weather conditions, and real operator
feedback. Such validation would help measure alert usefulness, identify false
positive causes, and determine how the dashboard supports decision-making in
real monitoring situations.
